\documentclass[12pt]{article}
\usepackage[utf8]{inputenc}
\usepackage[T1]{fontenc}
\usepackage{hyperref}
\title{Gizrra Language Resources for K-12 Curriculum Development}
\author{Compiled by Claude Sonnet 5 (claude-sonnet-5) for Andr\'as Kornai}
\date{}
\begin{document}
\maketitle

\section{Introduction}
Gizrra (also called Toga; ISO 639-3: tof; Glottocode: gizr1240) is an Eastern Trans-Fly language of Western Province, Papua New Guinea, south of the Fly River. It has about 1,100 speakers and two varieties, Western Gizrra and Waidoro. The main villages are Kulalae, Waidorro and Kolmet. It has a Latin-based orthography. See \url{https://en.wikipedia.org/wiki/Gizrra_language}.

\section{Primary Educational Resources}
The search found no primers or textbooks, despite the prompt's note that the language is used in instruction.

\section{Linguistic and Reference Materials}
A Gizrra-English Dictionary (Nico and Elly van Bodegraven, 2021) is reported to exist as a PDF, but I did not find its URL. Glottolog lists a 1976 work by Fleischmann and Turpeinen: \url{https://glottolog.org/resource/reference/id/133002}. The Endangered Languages Project has an entry: \url{https://endangeredlanguages.com/lang/7369}.

\section{Religious and Community Resources}
Joshua Project has a people-group profile: \url{https://joshuaproject.net/people_groups/11906}.

\section{Digital and Audio Resources}
None found.

\section{Curriculum Development Context}
There is a small, stable community and a Latin orthography, and a dictionary reportedly exists. Teaching materials were not found and need a follow-up search.
\end{document}